\documentclass[11pt]{article}
\usepackage[T1]{fontenc}
\usepackage[utf8]{inputenc}
\usepackage{lmodern}
\usepackage[letterpaper,margin=1in]{geometry}
\usepackage{amsmath,amssymb,amsthm}
\usepackage{booktabs}
\usepackage{enumitem}
\usepackage{microtype}
\usepackage[compact]{titlesec}
\usepackage[hidelinks]{hyperref}

\setlist{itemsep=2pt,topsep=4pt}

\theoremstyle{plain}
\newtheorem{theorem}{Theorem}[section]
\newtheorem{lemma}[theorem]{Lemma}
\newtheorem{proposition}[theorem]{Proposition}
\newtheorem{corollary}[theorem]{Corollary}
\theoremstyle{definition}
\newtheorem{definition}[theorem]{Definition}
\theoremstyle{remark}
\newtheorem{remark}[theorem]{Remark}

\newcommand{\PP}{\mathcal P}
\newcommand{\Del}{\Delta_{\mathcal P}}
\newcommand{\bG}[1]{\overline{G}_{#1}}
\newcommand{\lk}{\operatorname{lk}}
\newcommand{\dl}{\operatorname{dl}}
\newcommand{\Orb}{\operatorname{Orb}}
\newcommand{\F}{\mathbb F}
\newcommand{\file}[1]{\path{#1}}

\title{Non-evasive monotone graph properties\\on ten vertices}
\author{Lee Marshall}
\date{Draft of 7 November 2026}

\begin{document}
\maketitle

\begin{abstract}
\noindent
A graph property on $n$ vertices is \emph{evasive} if every decision tree for it must, in the worst case, ask about all $\binom n2$ possible edges.
Karp's conjecture says that every nontrivial monotone graph property is evasive.
It is known when $n$ is a prime power and when $n=6$.
The smallest open case is $n=10$.
Angel and Borja showed that a non-evasive monotone property on ten vertices has one of nine types, and excluded four of them.
We exclude four more, which leaves one type, called type~5.
For type~5 we show the following.
Such a property contains every disconnected graph and no graph with disconnected complement.
It contains nine explicit graphs and none of their complements.
Its trace on a family of $32$ graph classes is one of eight patterns.
It is not self-dual.
Apart from the cited fixed point theorems, every step is a finite computation with explicit permutation groups.
A Python program of about 1,200 lines, which uses only the standard library, replays the certificates.
We do not decide type~5.
\end{abstract}

\section{Introduction}

A property of graphs on $n$ labelled vertices is \emph{monotone} if it survives the removal of edges (or, in the dual version, the addition of edges).
It is \emph{nontrivial} if some graphs have it and some do not.
A decision tree for a property asks, one edge at a time, whether the edge is present, and must in the end say whether the graph has the property.
The property is \emph{evasive} if every decision tree has depth $\binom n2$.
Karp's evasiveness conjecture says that every nontrivial monotone graph property is evasive.
Miller's survey \cite{Mil} gives the history.

Rivest and Vuillemin \cite{RV} proved that a nontrivial monotone property needs a constant fraction of the $\binom n2$ queries.
Kahn, Saks and Sturtevant \cite{KSS} showed that a non-evasive monotone property gives a collapsible simplicial complex.
With fixed point theorems of Smith and Oliver they proved the conjecture for prime powers $n$ and for $n=6$.
The smallest case that is still open is $n=10$.

Two papers treat $n=10$.
Engstr\"om \cite{Eng} applies Oliver's theorems to several explicit groups.
He shows that the indicator vector of the $22$ vertex-transitive graphs on ten vertices must be one of six vectors (his Theorem~14).
Angel and Borja \cite{AB} use a Sylow $5$-subgroup of $S_{10}$ and the Euler characteristic.
They show that a non-evasive property meets ten particular graph classes in one of nine ways, which we call types $1$ to $9$.
They exclude types $1,3,7,9$ with Oliver groups and Smith's theorem.
They write that they did not find groups that exclude types $2,4,5,6,8$ \cite[Remark~5.2]{AB}.
The same remark is in Borja's thesis \cite[Remark~1.8.2]{Bor}.

This note proves the following (precise statements are in Theorem~\ref{thm:main}).
Types $2,4,6,8$ cannot occur.
The groups have order $112$ and $192$, and the argument can be checked by hand (Section~\ref{sec:t2468}).
So a non-evasive property has type~5.
Such a property contains every disconnected graph and no graph whose complement is disconnected.
It contains nine explicit graphs, among them $C_{10}$, the Petersen graph and $K_5\,\square\,K_2$, and none of their complements.
Its trace on $32$ graph classes is one of eight patterns (Appendix~\ref{app:patterns}).
It is not self-dual.
We also re-derive the exclusion of types $1,3,7,9$ with replayable certificates.

\medskip\noindent
\textbf{Credit.}
The reduction to nine types, the Sylow $5$-subgroup argument and the exclusion of types $1,3,7,9$ are due to Angel and Borja.
The use of Oliver's theorem on the vertex-transitive graphs of ten vertices is due to Engstr\"om.
We use their method.
Angel and Borja use the link of one fixed vertex of the complex once, for type~1.
Our additions are the link and deletion equations for every group that fixes an edge, a computer search for contradicting groups among the conjugacy classes of subgroups of $S_{10}$ (GAP lists $1{,}593$; this count is GAP output and is not replayed), and the results listed above.

\medskip\noindent
\textbf{What is not claimed.}
We do not decide type~5, and we do not know whether a type-5 property exists.
Section~\ref{sec:open} says what has been tried.
The note has not been peer reviewed.
Each result is reduced to a finite check on explicit permutation groups, and a Python program replays the checks (Section~\ref{sec:replay}).
The cited theorems of Smith and Oliver are not re-proved.

\section{Setting and tools}

\subsection{Properties and complexes}

Let $V=\{0,1,\dots,9\}$ and let $E$ be the set of the $45$ two-element subsets of $V$.
A \emph{graph} is a subset of $E$.
The symmetric group $S_{10}$ acts on $V$, on $E$ and on graphs.
A \emph{property} is a set $\PP$ of graphs.
All properties in this note are
invariant under $S_{10}$,
closed under taking subgraphs,
nonempty, and
not equal to the set of all graphs.
We call these \emph{nontrivial monotone}.
Such a property contains the empty graph and does not contain $K_{10}$.
A monotone increasing property is covered by passing to the set of graphs that are not in it.
This does not change the depth of decision trees.

A \emph{decision tree} for $\PP$ asks whether an edge is present, and ends in a leaf marked ``in $\PP$'' or ``not in $\PP$''.
The property is \emph{non-evasive} if some decision tree for it has depth at most $44$, and \emph{evasive} otherwise.
Karp's conjecture for $n=10$ says that no nontrivial monotone property is non-evasive.

The simplicial complex $\Del$ has vertex set $E$.
Its faces are the graphs in $\PP$.
The group $S_{10}$ acts on it.
We write $\chi(K)=\sum_{\sigma\ne\emptyset}(-1)^{\dim\sigma}$, with the empty face not counted.
For an edge $e$ the \emph{link} and \emph{deletion} are
\[
\lk(e)=\{\sigma\in\Del:\ e\notin\sigma,\ \sigma\cup\{e\}\in\Del\},\qquad
\dl(e)=\{\sigma\in\Del:\ e\notin\sigma\}.
\]

\subsection{Non-evasive implies acyclic}

\begin{lemma}[Kahn--Saks--Sturtevant \cite{KSS}]\label{lem:kss}
Let $W$ be a finite set and $f:2^W\to\{0,1\}$ a function with $f(\emptyset)=1$ and $f(S)\ge f(T)$ whenever $S\subseteq T$.
Let $K_f=\{S: f(S)=1\}$.
If some decision tree for $f$ has depth less than $|W|$, then $K_f$ is collapsible.
In particular $K_f$ is contractible, acyclic over every coefficient ring, and $\chi(K_f)=1$.
\end{lemma}

\begin{proof}
Induction on $|W|$.
If $|W|=1$, the tree has depth $0$, so $f$ is constant, hence $f\equiv1$ and $K_f$ is a point.
Let $|W|>1$ and let $w$ be the variable asked at the root.
Put $f_0=f|_{x_w=0}$ and $f_1=f|_{x_w=1}$.
Both are decreasing, have trees of depth less than $|W|-1$, and $f_0(\emptyset)=1$.
Now $K_{f_0}=\dl(w)$ is collapsible by induction.
If $f_1\equiv0$, then $K_f=K_{f_0}$.
Otherwise $f_1(\emptyset)=1$ and $K_{f_1}=\lk(w)$ is collapsible by induction.
Fix a sequence of elementary collapses of $\lk(w)$ to a point $p$.
Each step removes a pair $\tau\subset\sigma$ with $\tau$ free in $\sigma$.
In $K_f$ the same steps with $w$ added, that is, removing $\tau\cup\{w\}\subset\sigma\cup\{w\}$, are also elementary collapses.
They remove every face containing $w$ except $\{w\}$ and $\{w,p\}$.
Then $\{w\}$ is free in $\{w,p\}$, and removing that pair leaves $\dl(w)$.
So $K_f$ collapses onto $\dl(w)$, which is collapsible.
\end{proof}

\begin{lemma}\label{lem:linkdel}
Let $\PP$ be non-evasive and nontrivial monotone.
Then $\Del$, and for every edge $e$ also $\lk(e)$ and $\dl(e)$, are collapsible.
\end{lemma}

\begin{proof}
Apply Lemma~\ref{lem:kss} to $\Del$ with $W=E$.
For $\lk(e)$ and $\dl(e)$ take a tree $T$ of depth at most $44$.
Its root asks about some edge $e_0$.
Choose $\pi\in S_{10}$ with $\pi(e_0)=e$.
Relabelling $T$ by $\pi$ gives a tree for $\pi\PP=\PP$ whose root asks about $e$.
The two subtrees below the root have depth at most $43$.
They decide membership in $\dl(e)$ and in $\lk(e)$, as functions on subsets of $E\setminus\{e\}$.
For $\lk(e)$ we need $f(\{e\})=1$.
If a one-edge graph were not in $\PP$, then $\PP$ would contain only the empty graph.
That property is evasive, since the all-zero input forces all $45$ queries.
\end{proof}

\subsection{Fixed point theorems}

\begin{theorem}[Smith \cite{Smi}]\label{thm:smith}
Let a finite $p$-group $Q$ act simplicially on a finite complex $K$ that is acyclic over $\F_p$.
Then the fixed point set $|K|^Q$ is acyclic over $\F_p$.
\end{theorem}

\begin{theorem}[Oliver \cite{Oli}; see {\cite[Thm.~4.1]{AB}}, {\cite[Thm.~4]{Eng}}]\label{thm:oliver}
Let $\Gamma$ be a finite group with a normal $p$-subgroup $N$ such that $\Gamma/N$ is cyclic.
Let $\Gamma$ act simplicially on a finite complex $K$ that is acyclic over $\F_p$.
Then $\chi(|K|^\Gamma)=1$.
\end{theorem}

We call a group as in Theorem~\ref{thm:oliver} an \emph{Oliver group}.
Every $p$-group is one, and so is every cyclic group.
A certificate names a prime $p$, the normal subgroup $N$, and an element $g$ with $\Gamma=\langle g\rangle N$.

\subsection{The equations}

Let $\Gamma\le S_{10}$.
Write $\Orb(\Gamma)$ for the set of $\Gamma$-orbits on $E$, and $\bigcup T$ for the graph that is the union of the orbits in a set $T$.
A point of $|K|$ fixed by $\Gamma$ has barycentric coordinates that are constant on $\Gamma$-orbits of vertices.
Its support is a union of orbits that is a face.
Adding the coordinates over each orbit identifies $|K|^\Gamma$ with the realisation of the complex whose vertices are the orbits that are faces and whose faces are the sets of such orbits with union a face \cite[\S2]{AB}.
For $K=\Del$ this gives the following.

\begin{proposition}\label{prop:eq}
Let $\PP$ be non-evasive and nontrivial monotone, and let $\Gamma$ be an Oliver group.
Then
\begin{equation}\tag{fix}
\sum_{\emptyset\ne T\subseteq\Orb(\Gamma),\ \bigcup T\in\PP}(-1)^{|T|-1}=1 .
\end{equation}
If $\Gamma$ fixes an edge $e$, put $\Orb'=\Orb(\Gamma)\setminus\{\{e\}\}$.
Then
\begin{align}\tag{lk}
\sum_{\emptyset\ne T\subseteq\Orb',\ \{e\}\cup\bigcup T\in\PP}(-1)^{|T|-1}&=1,\\
\tag{dl}
\sum_{\emptyset\ne T\subseteq\Orb',\ \bigcup T\in\PP}(-1)^{|T|-1}&=1 .
\end{align}
\end{proposition}

\begin{proof}
By Lemma~\ref{lem:linkdel} the complexes $\Del$, $\lk(e)$, $\dl(e)$ are acyclic over every field.
If $\Gamma$ fixes $e$, it acts on $\lk(e)$ and $\dl(e)$.
Apply Theorem~\ref{thm:oliver} and the description of $|K|^\Gamma$ above.
\end{proof}

Each equation is linear in the numbers $x_C=[C\in\PP]\in\{0,1\}$, where $C$ runs over isomorphism classes of graphs of the form $\bigcup T$ (or $\{e\}\cup\bigcup T$).
The coefficients are signed counts.
Monotonicity adds the inequalities $x_C\le x_{C'}$ when $C'$ is a subgraph of $C$.

\subsection{Duality}

For a property $\PP$ let $\PP^*=\{G:\ \overline G\notin\PP\}$, where $\overline G=E\setminus G$.

\begin{lemma}\label{lem:dual}
If $\PP$ is nontrivial monotone, so is $\PP^*$, and $(\PP^*)^*=\PP$.
Moreover $\PP^*$ is non-evasive if and only if $\PP$ is.
\end{lemma}

\begin{proof}
If $G'\subseteq G$ then $\overline{G'}\supseteq\overline G$, so $\overline{G}\notin\PP$ implies $\overline{G'}\notin\PP$.
Hence $\PP^*$ is closed under subgraphs.
The empty graph is in $\PP^*$ because $K_{10}\notin\PP$, and $K_{10}\notin\PP^*$ because the empty graph is in $\PP$.
Given a tree for $\PP$, keep its shape, answer each query about $e$ by the opposite answer, and flip the leaf labels.
This is a tree of the same depth for $\PP^*$.
\end{proof}

\section{The nine types of Angel and Borja}\label{sec:types}

Let $U=\{0,\dots,4\}$ and $U'=\{5,\dots,9\}$.
Let $G_1,\dots,G_5$ be the following graphs on $V$ (disjoint unions, with isolated vertices where needed):
\[
G_1=C_5\cup5K_1,\quad G_2=K_5\cup5K_1,\quad G_3=C_5\cup C_5,\quad G_4=C_5\cup K_5,\quad G_5=K_5\cup K_5 .
\]
Let $\overline{G}_i$ be the complement of $G_i$; for example $\overline G_5=K_{5,5}$.
Let $\mathcal G=\{G_1,\dots,G_5,\overline G_1,\dots,\overline G_5\}$, taken up to isomorphism.

Let $S=\langle(0\,1\,2\,3\,4),(5\,6\,7\,8\,9)\rangle$, a Sylow $5$-subgroup of $S_{10}$.
It has five orbits on $E$: the two $5$-cycles $c_1=\{i,i+1\}$ and $c_2=\{i,i+2\}$ inside $U$, the two $5$-cycles $c_1',c_2'$ inside $U'$, and the $25$ edges $m$ between $U$ and $U'$.
Every nonempty union of orbits is isomorphic to $K_{10}$ or to a member of $\mathcal G$.
For $G\in\mathcal G$ let $w(G)=\sum(-1)^{|T|-1}$, summed over the sets $T$ of orbits with $\bigcup T\cong G$.
Counting the sets $T$ gives
\[
w(G_1)=4,\quad w(G_2)=-2,\quad w(G_3)=-4,\quad w(G_4)=4,\quad w(G_5)=-1 .
\]
For example, $G_3=C_5\cup C_5$ comes from the four sets $\{c_i,c'_j\}$, each with sign $-1$, and $G_4$ comes from the four sets $\{c_i,c'_1,c'_2\}$ and $\{c_1,c_2,c'_j\}$, each with sign $+1$.
Since $\overline{\bigcup T}=\bigcup(\Orb(S)\setminus T)$, we also have $w(\overline G_i)=-w(G_i)$.

Let $\PP$ be non-evasive and nontrivial monotone.
The group $S$ is a $5$-group, so (fix) holds for $S$.
The empty union is not counted, and $K_{10}\notin\PP$.
Hence
\begin{equation}\label{eq:sylow}
\sum_{G\in\PP\cap\mathcal G}w(G)=1 .
\end{equation}
Since $\PP$ is closed under subgraphs, $\PP\cap\mathcal G$ is closed downward in the subgraph order on $\mathcal G$.
There are exactly nine downward closed subsets of $\mathcal G$ with weight sum $1$.
They are listed in Table~\ref{tab:types}, in the numbering of \cite[\S5]{AB}.
We say $\PP$ has \emph{type $k$} if $\PP\cap\mathcal G$ is the $k$-th set.
We recomputed the weights, the order and the nine sets by program, and the result agrees with \cite{AB}.

\begin{table}[t]
\centering
\begin{tabular}{cll}
\toprule
type & $\PP\cap\mathcal G$ & status\\
\midrule
1 & $\overline G_5$ & excluded in \cite{AB}\\
2 & $G_1,\ \overline G_4,\ \overline G_5$ & excluded here\\
3 & $G_1,\ G_3,\ \overline G_5$ & excluded in \cite{AB}\\
4 & $G_1,G_2,\ \overline G_2,\overline G_4,\overline G_5$ & excluded here\\
5 & $G_1,G_2,G_3,G_4,G_5$ & open\\
6 & $G_1,G_3,\ \overline G_3,\overline G_4,\overline G_5$ & excluded here\\
7 & $G_1,G_2,G_3,G_4,\ \overline G_2,\overline G_4,\overline G_5$ & excluded in \cite{AB}\\
8 & $G_1,G_2,G_3,\ \overline G_2,\overline G_3,\overline G_4,\overline G_5$ & excluded here\\
9 & $G_1,G_2,G_3,G_4,\ \overline G_1,\overline G_2,\overline G_3,\overline G_4,\overline G_5$ & excluded in \cite{AB}\\
\bottomrule
\end{tabular}
\caption{The nine types.}\label{tab:types}
\end{table}

\begin{lemma}\label{lem:typedual}
If $\PP$ has type $k$, then $\PP^*$ has type $10-k$.
In particular, if $\PP$ has type~$5$ then so does $\PP^*$.
\end{lemma}

\begin{proof}
We have $G_i\in\PP^*$ iff $\overline G_i\notin\PP$, and $\overline G_i\in\PP^*$ iff $G_i\notin\PP$.
Apply this to the nine sets in Table~\ref{tab:types}.
For example, type $2$ is $\{G_1,\overline G_4,\overline G_5\}$, and the dual set is $\{G_1,G_2,G_3,\overline G_2,\overline G_3,\overline G_4,\overline G_5\}$, which is type $8$.
\end{proof}

\section{Results}

\begin{theorem}\label{thm:main}
Let $\PP$ be a non-evasive nontrivial monotone graph property on ten vertices.
Then the following hold.
\begin{enumerate}[label=(\alph*)]
\item $\PP$ has type~$5$.
\item $\PP$ contains every disconnected graph, and no graph whose complement is disconnected.
In particular $K_a\cup K_{10-a}\in\PP$ for $1\le a\le5$, and $\PP$ contains no graph with a vertex adjacent to all others.
\item $\PP$ contains each of the following nine graphs and none of their complements:
$C_{10}$; $C_5$ with a pendant edge at each vertex; the Petersen graph; the pentagonal prism $C_5\,\square\,K_2$; $K_5$ with a pendant edge at each vertex; $C_5$ with a triangle on each edge; the square of $C_{10}$ (the circulant with differences $\pm1,\pm2$); $(K_5\cup C_5)$ plus a perfect matching between the two parts; $K_5\,\square\,K_2$.
\item Let $Q=\langle(0\,1\,2\,3\,4)(5\,6\,7\,8\,9)\rangle$.
Among the $32$ isomorphism classes of connected graphs with connected complement that are unions of $Q$-orbits, the classes that lie in $\PP$ form one of eight explicit sets (Appendix~\ref{app:patterns}).
\item $\PP\ne\PP^*$.
\end{enumerate}
\end{theorem}

Parts (a) and (b) are proved in Sections~\ref{sec:t2468}--\ref{sec:t5}, parts (c) and (d) in Section~\ref{sec:z5}, and part (e) in Section~\ref{sec:nsd}.
The proofs depend on Theorems~\ref{thm:smith} and~\ref{thm:oliver}, and on the computer certificates.
So a counterexample to Karp's conjecture on ten vertices, if there is one, has the form $\PP=D\cup X$.
Here $D$ is the set of disconnected graphs and $X$ is a downward closed set of connected graphs with connected complement.
It satisfies (c)--(e).

\section{Proofs}

\subsection{Certificates}\label{sec:cert}

A certificate for type $T$ lists groups $\Gamma\le S_{10}$ by generators, each with the data of an Oliver group (Theorem~\ref{thm:oliver}), and a list of equations (fix), (lk) or (dl), the last two with the fixed edge.
The program in \file{replay/replay_cert.py} works as follows.
(i)~It recomputes the nine types of Section~\ref{sec:types}.
(ii)~For each group it closes the generators by brute force, checks the order, checks that $N$ is normal of order $p^k$ with $p$ prime and that $\Gamma=\langle g\rangle N$, and computes the edge orbits.
(iii)~For each equation it forms the graph of every set of orbits and sorts these graphs into isomorphism classes (an embedding found by backtracking, with equal edge counts).
(iv)~It gives each class a status for the type $T$: \emph{in} if the class embeds in a member of the type, \emph{out} if a graph of $\mathcal G\cup\{K_{10}\}$ outside the type embeds in it, and \emph{free} otherwise.
A witness permutation is checked in the first two cases.
(v)~Each equation becomes a constant plus a signed sum of free $0/1$ unknowns, equal to $1$.
(vi)~A complete search with propagation tries all $0/1$ values of the free unknowns, subject to the equations and to $x_C\le x_{C'}$ among free classes.
If there is no assignment, the type is excluded.
Leaving out the inequalities between free and fixed classes can only add assignments, so ``no assignment'' remains a valid conclusion.
A certificate for a statement ``$H\in\PP$'' adds the assumption $H\notin\PP$.
Table~\ref{tab:cert} lists the certificates for types $1$--$4$ and $6$--$9$.
The group names are positions in the list of conjugacy classes of subgroups of $S_{10}$ made by \file{gap/enum_subgroups.g} (GAP 4.12.1); the replay does not use them.

\begin{table}[t]
\centering
\begin{tabular}{lllll}
\toprule
types & group & order & equation & source\\
\midrule
1, 9 & $\langle(0\,0'),(1\,2\,3\,4\,1'\,2'\,3'\,4')\rangle\cong C_2\times C_8$ & 16 & (lk) & \cite{AB}\\
1, 9 & class 97 & 8 & (lk) & this work\\
3, 7 & $\langle\sigma,\rho,\rho'\rangle$ (Section~\ref{sec:t1379}) & 50 & (fix) & \cite{AB}\\
3, 7 & class 836 & 50 & (fix) & this work\\
2, 8 & classes 1132 and 1279 & 112, 192 & (lk) for each & this work\\
4, 6 & class 1279 & 192 & (lk) & this work\\
4, 6 & an order-16 group found by a computer search & 16 & (lk) & this work\\
\bottomrule
\end{tabular}
\caption{Certificates for the excluded types (files \texttt{certs/type*.json}); (lk) is taken at a fixed edge.}\label{tab:cert}
\end{table}

\subsection{Types 1, 3, 7, 9}\label{sec:t1379}

These are due to Angel and Borja.
We give the argument for types $3$ and $7$.
Write $0',\dots,4'$ for $5,\dots,9$.
Let $\sigma=(0\,0')(1\,1')(2\,2')(3\,3')(4\,4')$, $\rho=(0\,1\,2\,3\,4)$, $\rho'=(0'\,1'\,2'\,3'\,4')$ and $\Gamma=\langle\sigma,\rho,\rho'\rangle$.
It has order $50$, and $N=C_5\times C_5$ is normal with quotient $C_2$.
Its edge orbits are $A=K_{5,5}$, $B$ (the $5$-cycles $\{i,i+1\}$ in $U$ and $U'$) and $C$ (the $5$-cycles $\{i,i+2\}$ in $U$ and $U'$).
Then $A\cong\overline G_5$, $B\cong C\cong G_3$, $A\cup B\cong A\cup C\cong\overline G_3$ and $B\cup C\cong G_5$.
Types $3$ and $7$ contain $G_3$ and $\overline G_5$, but neither $\overline G_3$ nor $G_5$.
So the fixed complex is three isolated points.
Its Euler characteristic is $3$, not $1$.
For types $1$ and $9$ the argument of \cite{AB} uses a group of order $16$ and the link of the fixed edge; the replay checks it, and also a group of order $8$.
Type $9$ is dual to type $1$ and type $7$ to type $3$ (Lemma~\ref{lem:typedual}); the replay also checks the dual types directly.

\subsection{Types 2, 4, 6, 8}\label{sec:t2468}

Let the vertices be two \emph{hubs} $a,b$ and a set $W$ of eight vertices.
Let $e=\{a,b\}$ and let $\tau=(a\,b)$.

\smallskip\noindent
\textbf{The group $\Gamma_1$} (order $112$).
Identify $W$ with $\F_8$.
Let $\Gamma_1=\mathrm{AGL}(1,8)\times\langle\tau\rangle$, where $\mathrm{AGL}(1,8)=\{x\mapsto\alpha x+\beta\}$ acts on $W$.
It is $2$-transitive on $W$.
The edge orbits are $\{e\}$, $A$ (the $16$ edges from hubs to $W$) and $B$ (the $28$ edges inside $W$).
The subgroup $N_1=\{x\mapsto x+\beta\}\times\langle\tau\rangle$ is normal of order $16$, and $\Gamma_1/N_1\cong\F_8^\times\cong C_7$.

\smallskip\noindent
\textbf{The group $\Gamma_2$} (order $192$).
Identify $W$ with $\F_2\times\F_4$.
Let $\Gamma_2'$ be the set of maps $(t,z)\mapsto(t+s,\ \lambda z+\mu t+\nu)$ with $s\in\F_2$, $\lambda\in\F_4^\times$, $\mu,\nu\in\F_4$.
This is a group of order $96$.
Let $\Gamma_2=\Gamma_2'\times\langle\tau\rangle$.
The subgroup $N_2$ with $\lambda=1$ is normal of order $64$, and $\Gamma_2/N_2\cong C_3$.
The group $\Gamma_2'$ preserves the partition of $W$ into two sets of four, $\{t=0\}$ and $\{t=1\}$.
It is transitive on the pairs inside a block and on the pairs across blocks.
The edge orbits are $\{e\}$, $A$ as before, $B'=K_{4,4}$ (the $16$ edges across the blocks) and $C'=2K_4$ (the $12$ edges inside the blocks).

Both groups are Oliver groups for $p=2$ and fix $e$.
The script \file{note/check_explicit.py} builds both from these definitions, checks the claims above, and replays the type-2, 4, 6, 8 equations for them.
The certificates in \file{certs/} use conjugates found by GAP.

\begin{proposition}\label{prop:t24}
There is no non-evasive nontrivial monotone property of type $2$, $4$, $6$ or $8$.
\end{proposition}

\begin{proof}
By Lemma~\ref{lem:dual} and Lemma~\ref{lem:typedual} it is enough to treat types $2$ and $4$.
Fix a non-evasive $\PP$.
We use the link equation (lk) for $\Gamma_1$ and $\Gamma_2$ and the edge $e$.
The complex $\lk(e)^{\Gamma}$ has one vertex for each orbit $O\ne\{e\}$ with $\{e\}\cup O\in\PP$.
Its faces are the sets of such orbits whose union with $e$ lies in $\PP$.
The graphs involved are these.
\begin{itemize}[nosep]
\item $e\cup A=K_{1,1,8}$, the join of an edge with eight independent vertices.
\item $e\cup B'=K_2\cup K_{4,4}$.
\item $e\cup C'=K_2\cup2K_4$ and $e\cup B=K_2\cup K_8$.
\item $e\cup A\cup B'$ contains $2C_5=G_3$.
Take the sides $X,Y$ of $K_{4,4}$ and cycles $(a,x_1,y_1,x_2,y_2)$ and $(b,x_3,y_3,x_4,y_4)$.
\item $e\cup A\cup C'$ contains $2K_5=G_5$: the hub $a$ with one $K_4$, the hub $b$ with the other.
\item $e\cup B'\cup C'=K_2\cup K_8$, which contains $K_5\cup5K_1=G_2$.
\end{itemize}

\emph{Type 2} ($\PP\cap\mathcal G=\{G_1,\overline G_4,\overline G_5\}$).
For $\Gamma_1$: $e\cup B\supseteq G_2\notin\PP$, so $B$ is not a vertex.
So $\lk(e)^{\Gamma_1}$ is empty or the single vertex $A$, and $\chi=1$ forces $K_{1,1,8}\in\PP$.
For $\Gamma_2$: $A$ is a vertex by what we just proved.
$B'$ is a vertex, because $K_2\cup K_{4,4}\subseteq K_{5,5}=\overline G_5\in\PP$.
No two of $A,B',C'$ span an edge: the unions contain $G_3$, $G_5$ and $G_2$, which are not in $\PP$.
So $\lk(e)^{\Gamma_2}$ has at least two vertices and no edges.
Then $\chi\ge2$, a contradiction.

\emph{Type 4} ($\PP\cap\mathcal G=\{G_1,G_2,\overline G_2,\overline G_4,\overline G_5\}$).
Use $\Gamma_2$ only.
The graph $\overline G_2$ is the join of $K_5$ with five independent vertices, and $K_{1,1,8}\subseteq\overline G_2$ (send the hubs to two vertices of the $K_5$).
So $A$ is a vertex.
$B'$ is a vertex, because $K_2\cup K_{4,4}\subseteq K_{5,5}=\overline G_5\in\PP$.
The unions $A\cup B'$ and $A\cup C'$ contain $G_3$ and $G_5$, which are not in $\PP$, so they are not faces.
The only possible edge is $B'C'$, which is a face iff $y=[K_2\cup K_8\in\PP]=1$.
Put $x=[K_2\cup2K_4\in\PP]$.
Then $C'$ is a vertex iff $x=1$, and $y\le x$ because $K_2\cup2K_4\subseteq K_2\cup K_8$.
There is no triangle.
So $\chi=2+x-y\ge2$, a contradiction.
\end{proof}

\subsection{Type 5: the sandwich}\label{sec:t5}

\begin{proposition}\label{prop:A}
Let $\PP$ be non-evasive, nontrivial, monotone, and of type~$5$.
Then $K_a\cup K_{10-a}\in\PP$ for $a=1,2,3,4$.
\end{proposition}

\begin{proof}
For each $a$ there is a certificate \file{certs/type5_force_K<a>_K<10-a>.json}.
It consists of $6$ to $9$ equations of the kinds (fix), (lk), (dl) for $5$ to $8$ Oliver groups of orders $8$ to $192$ (primes $p=2,3,5,7$), together with the assumption $K_a\cup K_{10-a}\notin\PP$.
The replay finds no solution; the final search runs over $68$ to $80$ free classes.
The replay checks by an isomorphism test that the assumed graph is $K_a\cup K_{10-a}$, with $a$ taken from the command line of the manifest and not from the certificate.
\end{proof}

\begin{proof}[Proof of Theorem~\ref{thm:main}(a) and (b)]
By Section~\ref{sec:types}, $\PP$ has one of the types $1,\dots,9$. Types $1,3,7,9$ are excluded in Section~\ref{sec:t1379}, and types $2,4,6,8$ by Proposition~\ref{prop:t24}. This leaves type~$5$, which is (a).
For (b), let $H$ be a disconnected graph.
Group its components into two sets with $a$ and $10-a$ vertices, $a\le5$.
Then $H$ is a subgraph of a copy of $K_a\cup K_{10-a}$, which is in $\PP$ by Proposition~\ref{prop:A} (for $a\le4$) and by type~5 (for $a=5$, since $G_5=2K_5$).
So $H\in\PP$.
Now $\PP^*$ is non-evasive, nontrivial, monotone (Lemma~\ref{lem:dual}) and of type~5 (Lemma~\ref{lem:typedual}).
So every disconnected $H$ lies in $\PP^*$, that is, $\overline H\notin\PP$.
Every graph with disconnected complement has this form.
\end{proof}

Write $\mathrm{CC}$ for the set of connected graphs with connected complement.
By (b), $\PP=D\cup X$ with $X=\PP\cap\mathrm{CC}$ downward closed in $\mathrm{CC}$.

\subsection{Type 5: the group of order 5}\label{sec:z5}

Let $q=(0\,1\,2\,3\,4)(5\,6\,7\,8\,9)$ and $Q=\langle q\rangle$.
$Q$ acts freely on $E$, so it has nine orbits of five edges.
Of the $511$ nonempty unions of orbits, $20$ are disconnected, $21$ are connected with disconnected complement, and $470$ lie in $\mathrm{CC}$.
The $470$ fall into $32$ isomorphism classes.
By Lemma~\ref{lem:linkdel} and Theorem~\ref{thm:smith} (with $p=5$) the fixed complex $\Del^Q$ is acyclic over $\F_5$.
By (b) this complex is determined by the set $X_Q$ of the $32$ classes that lie in $\PP$: its faces are the sets of orbits whose union is disconnected, together with those whose union is in a class of $X_Q$.
The set $X_Q$ is closed downward under embedding.
The subcomplex of disconnected unions has reduced $\F_5$-homology of rank $5$ in dimension $0$, and the relative complex of the other faces must account for it.

\begin{proof}[Proof of Theorem~\ref{thm:main}(c) and (d)]
The program lists all $2{,}539$ downward closed subsets of the $32$ classes.
For each it computes the $\F_5$-homology of the relative complex (the faces whose union is in $\mathrm{CC}$) and keeps the subset if this equals the homology of the disconnected part, shifted by one.
By the long exact sequence of the pair, every subset with an acyclic fixed complex passes this test.
The test is only necessary: equal dimensions do not force the connecting maps to be isomorphisms.
So its survivors may include subsets that are not admissible, which does no harm to a statement that holds for all of them.
Sixteen subsets pass.
The script \file{replay/check_z5_absolute.py} computes the $\F_5$-homology of the whole fixed complex for each of the $2{,}539$ subsets.
It finds that exactly the same $16$ are acyclic.
So for this complex the necessary test is also sufficient, and this is replayed.
Two further Oliver groups, of order $80$ and $160$ with three edge orbits each, give equations (fix) whose unions are all among the $32$ classes.
Eight of the $16$ patterns satisfy both.
The program checks that each of the nine graphs of (c) lies in all eight patterns and that each of their complements lies in none.
The nine graphs are built from their definitions, and matched with the edge lists of the certificate by isomorphism tests, in \file{replay/check_nine_graphs.py}.
The two certificates are \file{certs/lemma_z5.json} (seven graphs, using only Smith's theorem) and \file{certs/lemma_z5_euler.json} (nine graphs).
\end{proof}

The eight patterns have $9,11,11,13,19,21,21,23$ classes.
They form four pairs under $\PP\leftrightarrow\PP^*$.
They are listed in Appendix~\ref{app:patterns}.

\subsection{Type 5 is not self-dual}\label{sec:nsd}

\begin{proof}[Proof of Theorem~\ref{thm:main}(e)]
Taking complements maps the union of a set $T$ of $Q$-orbits to the union of the complementary set.
So it permutes the $32$ classes.
No graph on ten vertices is isomorphic to its complement, since that would need $45/2$ edges.
So the permutation has no fixed class, and it splits the $32$ classes into $16$ pairs.
Suppose $\PP=\PP^*$.
For every class $c$ we then have $c\in\PP$ iff $c\in\PP^*$ iff $\overline c\notin\PP$.
So $X_Q$ contains exactly one class of each pair, and $|X_Q|=16$.
By Theorem~\ref{thm:main}(d), $X_Q$ is one of the eight patterns, which have $9,11,11,13,19,21,21,23$ classes.
None has $16$.
\end{proof}

The program \file{replay/replay_nonselfdual.py} re-derives the patterns and checks this count.
It also checks that the patterns are closed under $X\mapsto\{c:\ \overline c\notin X\}$ and that none is fixed by it, which the argument above does not need.

\section{Replaying the proofs}\label{sec:replay}

Requirements: Python~3, standard library only (tested with Python~3.14.0).
Do not run it with \texttt{-O}; the checks are assertions, and \file{replay/replay_lib.py} refuses to start without them.
From the directory \file{karp10} of the archive:
\begin{verbatim}
python replay/replay_all.py
\end{verbatim}
It first checks that \file{certs/} and \file{replay/} contain exactly the files listed in \file{replay/RELEASE-MANIFEST.json} (the $20$ certificates and the scripts) and that each has the pinned SHA-256 hash.
There is no option that skips this check.
The checks are then run, in isolated interpreters, exactly as the manifest lists them.
The SHA-256 hash of the manifest file itself, which the program prints, is to be quoted wherever the archive is published.
Then it replays the $20$ certificates and three further checks (23 checks in all): non-self-duality (Section~\ref{sec:nsd}), absolute acyclicity of the $16$ patterns (Section~\ref{sec:z5}), and the identification of the nine graphs (Theorem~\ref{thm:main}(c)).
It takes between one and three minutes, depending on the load of the machine, prints one line per check, and ends with \texttt{ALL VALID}.
Single checks:
\begin{verbatim}
python replay/replay_cert.py certs/type4.json -v
python replay/replay_smith.py certs/lemma_z5_euler.json
python replay/replay_nonselfdual.py certs/lemma_z5_euler.json
python replay/check_z5_absolute.py
python replay/check_nine_graphs.py
python note/check_explicit.py
python note/z5_table.py
\end{verbatim}
The last two rebuild the groups of Section~\ref{sec:t2468} and the table of Appendix~\ref{app:patterns}.
The script \file{note/run_checks.sh} runs all of them.
Group closure, orbits, isomorphism tests and the search are written from scratch in \file{replay/replay_lib.py} and \file{replay/replay_cert.py}.
Nothing from GAP, nauty or a solver is used.
GAP and nauty were used only to find the groups and the certificates.
As a second opinion on the search, \file{note/crosscheck_z3.py} (it needs z3 and is not part of the replay) rebuilds the $0/1$ systems of the $18$ equation certificates with separate code and gives them to z3.
All $18$ are unsatisfiable, and the same systems against type~5 (or without the assumption in Proposition~\ref{prop:A}) are satisfiable.

\section{What remains open}\label{sec:open}

By Theorem~\ref{thm:main}, a counterexample on ten vertices is $\PP=D\cup X$ with $X$ a downward closed set of graphs in $\mathrm{CC}$, subject to (c)--(e).
We did not find a contradiction, and we did not find an example.

\medskip\noindent
\textbf{Computer evidence (not certificates).}
Nothing in this subsection is replayed by the programs of Section~\ref{sec:replay}.
It rests on output of an SMT solver and of search programs whose data files (about $41{,}000$ graph classes, tables made with GAP and nauty) are not part of the replay archive.
We tested further necessary conditions on $X$.
\begin{itemize}[nosep]
\item The equations (fix), (lk), (dl) for the $1{,}524$ conjugacy classes of subgroups of $S_{10}$ with at most $16$ edge orbits, and for four more classes (of orders $3$, $5$, $9$ and $10$).
For one pattern we also imposed (fix) for the involution with five $2$-cycles.
\item The relative-homology condition of Section~\ref{sec:z5}, which is necessary for acyclicity over $\F_p$, computed exactly for the fixed, link and deletion complexes of the $p$-groups among these groups ($1{,}647$ complexes) and for two more $p$-groups ($34$ complexes).
\item Monotonicity on the $41{,}134$ isomorphism classes of graphs in $\mathrm{CC}$ that occur as unions of orbits of these groups.
\end{itemize}
With an SMT solver we found $0/1$ assignments to these classes that satisfy all of this together with pattern $1$, $2$ or $3$ of Appendix~\ref{app:patterns}.
By duality the same holds for patterns $8$, $7$ and $6$.
For each of the patterns $1$ to $4$ (and so $5$ to $8$) we also found assignments for the system without the four extra classes and the two extra $p$-groups, which in addition satisfy two further families: the depth-2 equations of one selector branch, and the Euler equations of eight further $2$-groups.
The case of pattern~$4$ was found by a repair search that fixes most classes to a model satisfying one side of the system.
We re-checked each assignment with separate code.
These assignments are not graph properties known to be non-evasive.
They show only that the conditions above do not contradict each other.
The assignments fail three further tests, which we have not imposed: the same condition for the involution with five $2$-cycles, for a group of order $3$, and for eight $2$-groups of order $4$ or $8$.
Whether the assignments can be repaired is open.
The conditions from the transposition and from the trivial group are untested.
Conditions from the second query of a decision tree were tested only in part, and conditions from later queries not at all.
As a sanity check (also solver output, not replayed), the same kind of system, with all subgroup classes, has no solution for $n=5,6,7$; for $n=6$ it needs the link equations.

\medskip\noindent
\textbf{The obstruction.}
The complex of disconnected graphs on $n$ vertices is homotopy equivalent to a wedge of $(n-1)!$ spheres of dimension $n-3$ \cite{Vas,BBLSW}.
For $n=10$ this is $9!$ spheres of dimension $7$.
A type-5 property adds graphs of $\mathrm{CC}$ to this complex so that the result is acyclic and the $S_{10}$-symmetry is kept.
The conditions above are necessary for that.
A proof along these lines needs a contradiction among such conditions, or an argument of a different kind, for example an adversary strategy for properties between $D$ and its dual.

\section{Use of AI, and availability}

This work was carried out with AI assistance (Claude, Anthropic), directed by the author.
The certificates are replayable.
The note has not been peer reviewed.
The author is responsible for all content.
The code, the certificates and the logs are in the archive that accompanies this note: \texttt{[location to be added]}.

{\footnotesize
\begin{thebibliography}{99}
\setlength{\itemsep}{1pt}
\bibitem{AB} A.~Angel and J.~Borja, \emph{The evasiveness conjecture and graphs on $2p$ vertices}, J.~Graph Theory \textbf{91} (2019), 35--52; arXiv:1603.04412. Section and theorem numbers in this note refer to arXiv version~1.
\bibitem{BBLSW} E.~Babson, A.~Bj\"orner, S.~Linusson, J.~Shareshian and V.~Welker, \emph{Complexes of not $i$-connected graphs}, Topology \textbf{38} (1999), 271--299; arXiv:math/9705219.
\bibitem{Bor} J.~M.~Borja Soto, \emph{Equivariant algebraic topology applied to some problems in topological combinatorics}, PhD thesis, Universidad de los Andes, 2016.
\bibitem{Eng} A.~Engstr\"om, \emph{Transitive graphs in counterexamples to Karp's conjecture}, arXiv:math/0512421 (2005).
\bibitem{KSS} J.~Kahn, M.~Saks and D.~Sturtevant, \emph{A topological approach to evasiveness}, Combinatorica \textbf{4} (1984), 297--306.
\bibitem{Mil} C.~A.~Miller, \emph{Evasiveness of graph properties and topological fixed-point theorems}, Found.\ Trends Theor.\ Comput.\ Sci.\ \textbf{7} (2013), 337--415; arXiv:1306.0110.
\bibitem{Oli} R.~Oliver, \emph{Fixed-point sets of group actions on finite acyclic complexes}, Comment.\ Math.\ Helv.\ \textbf{50} (1975), 155--177.
\bibitem{RV} R.~L.~Rivest and J.~Vuillemin, \emph{On recognizing graph properties from adjacency matrices}, Theoret.\ Comput.\ Sci.\ \textbf{3} (1976), 371--384.
\bibitem{Smi} P.~A.~Smith, \emph{Fixed-point theorems for periodic transformations}, Amer.\ J.\ Math.\ \textbf{63} (1941), 1--8.
\bibitem{Vas} V.~A.~Vassiliev, \emph{Complexes of connected graphs}, in: The Gelfand Mathematical Seminars 1990--1992, Birkh\"auser, 1993, 223--235.
\end{thebibliography}
}

\appendix
\section{The 32 classes and the eight patterns}\label{app:patterns}

Table~\ref{tab:patterns} lists the $32$ isomorphism classes of connected graphs with connected complement that are unions of $Q$-orbits.
Here $e$ is the number of edges, $t$ the number of triangles, $q$ the number of $4$-cycles, and ``cpl.'' the class of the complement.
Degrees are written degree$^{\text{multiplicity}}$.
A star marks the nine graphs of Theorem~\ref{thm:main}(c).
Column $j$ marks the classes in pattern $j$.
Patterns are numbered by size, and pattern $j$ is dual to pattern $9-j$.
Pattern~1 is the nine graphs alone.
Patterns 2 and 3 add the classes $\{4,9\}$ and $\{4,15\}$, and pattern 4 adds $\{4,9,15,20\}$.
The file \file{note/z5_patterns.json} has an edge list for each class.

\begin{table}[!htb]
\centering\footnotesize\setlength{\tabcolsep}{3.5pt}\renewcommand{\arraystretch}{0.88}
\begin{tabular}{@{}rrlrrr@{\hspace{6pt}}cccccccc@{\hspace{14pt}}rrlrrr@{\hspace{6pt}}cccccccc@{}}
\toprule
class & $e$ & degrees & $t$ & $q$ & cpl. & 1&2&3&4&5&6&7&8 & class & $e$ & degrees & $t$ & $q$ & cpl. & 1&2&3&4&5&6&7&8 \\
\midrule
1$^{*}$ & 10 & $3^{5}\,1^{5}$ & 0 & 0 & 31 & $\bullet$ & $\bullet$ & $\bullet$ & $\bullet$ & $\bullet$ & $\bullet$ & $\bullet$ & $\bullet$ & 17 & 25 & $7^{5}\,3^{5}$ & 25 & 65 & 9 &  &  &  &  &  & $\bullet$ &  & $\bullet$ \\
2$^{*}$ & 10 & $2^{10}$ & 0 & 0 & 32 & $\bullet$ & $\bullet$ & $\bullet$ & $\bullet$ & $\bullet$ & $\bullet$ & $\bullet$ & $\bullet$ & 18 & 25 & $6^{5}\,4^{5}$ & 15 & 45 & 12 &  &  &  &  &  &  &  &  \\
3$^{*}$ & 15 & $5^{5}\,1^{5}$ & 10 & 15 & 25 & $\bullet$ & $\bullet$ & $\bullet$ & $\bullet$ & $\bullet$ & $\bullet$ & $\bullet$ & $\bullet$ & 19 & 25 & $6^{5}\,4^{5}$ & 15 & 55 & 11 &  &  &  &  & $\bullet$ & $\bullet$ & $\bullet$ & $\bullet$ \\
4 & 15 & $4^{5}\,2^{5}$ & 0 & 5 & 27 &  & $\bullet$ & $\bullet$ & $\bullet$ & $\bullet$ & $\bullet$ & $\bullet$ & $\bullet$ & 20 & 25 & $6^{5}\,4^{5}$ & 20 & 45 & 10 &  &  &  & $\bullet$ & $\bullet$ & $\bullet$ & $\bullet$ & $\bullet$ \\
5$^{*}$ & 15 & $4^{5}\,2^{5}$ & 5 & 0 & 26 & $\bullet$ & $\bullet$ & $\bullet$ & $\bullet$ & $\bullet$ & $\bullet$ & $\bullet$ & $\bullet$ & 21 & 25 & $5^{10}$ & 10 & 50 & 16 &  &  &  &  &  &  &  &  \\
6$^{*}$ & 15 & $3^{10}$ & 0 & 0 & 28 & $\bullet$ & $\bullet$ & $\bullet$ & $\bullet$ & $\bullet$ & $\bullet$ & $\bullet$ & $\bullet$ & 22 & 25 & $5^{10}$ & 15 & 35 & 15 &  &  &  &  &  &  & $\bullet$ & $\bullet$ \\
7$^{*}$ & 15 & $3^{10}$ & 0 & 5 & 30 & $\bullet$ & $\bullet$ & $\bullet$ & $\bullet$ & $\bullet$ & $\bullet$ & $\bullet$ & $\bullet$ & 23 & 25 & $5^{10}$ & 20 & 35 & 13 &  &  &  &  & $\bullet$ & $\bullet$ & $\bullet$ & $\bullet$ \\
8 & 15 & $3^{10}$ & 0 & 5 & 29 &  &  &  &  & $\bullet$ & $\bullet$ & $\bullet$ & $\bullet$ & 24$^{*}$ & 25 & $5^{10}$ & 20 & 40 & 14 & $\bullet$ & $\bullet$ & $\bullet$ & $\bullet$ & $\bullet$ & $\bullet$ & $\bullet$ & $\bullet$ \\
9 & 20 & $6^{5}\,2^{5}$ & 15 & 30 & 17 &  & $\bullet$ &  & $\bullet$ & $\bullet$ & $\bullet$ & $\bullet$ & $\bullet$ & 25 & 30 & $8^{5}\,4^{5}$ & 40 & 135 & 3 &  &  &  &  &  &  &  &  \\
10 & 20 & $5^{5}\,3^{5}$ & 5 & 15 & 20 &  &  &  &  &  & $\bullet$ & $\bullet$ & $\bullet$ & 26 & 30 & $7^{5}\,5^{5}$ & 30 & 115 & 5 &  &  &  &  &  &  &  &  \\
11 & 20 & $5^{5}\,3^{5}$ & 10 & 10 & 19 &  &  &  &  & $\bullet$ & $\bullet$ & $\bullet$ & $\bullet$ & 27 & 30 & $7^{5}\,5^{5}$ & 35 & 105 & 4 &  &  &  &  &  &  &  & $\bullet$ \\
12$^{*}$ & 20 & $5^{5}\,3^{5}$ & 10 & 20 & 18 & $\bullet$ & $\bullet$ & $\bullet$ & $\bullet$ & $\bullet$ & $\bullet$ & $\bullet$ & $\bullet$ & 28 & 30 & $6^{10}$ & 30 & 90 & 6 &  &  &  &  &  &  &  &  \\
13 & 20 & $4^{10}$ & 0 & 25 & 23 &  &  &  &  & $\bullet$ & $\bullet$ & $\bullet$ & $\bullet$ & 29 & 30 & $6^{10}$ & 30 & 95 & 8 &  &  &  &  & $\bullet$ & $\bullet$ & $\bullet$ & $\bullet$ \\
14 & 20 & $4^{10}$ & 0 & 30 & 24 &  &  &  &  &  &  &  &  & 30 & 30 & $6^{10}$ & 30 & 95 & 7 &  &  &  &  &  &  &  &  \\
15 & 20 & $4^{10}$ & 5 & 10 & 22 &  &  & $\bullet$ & $\bullet$ & $\bullet$ & $\bullet$ & $\bullet$ & $\bullet$ & 31 & 35 & $8^{5}\,6^{5}$ & 55 & 215 & 1 &  &  &  &  &  &  &  &  \\
16$^{*}$ & 20 & $4^{10}$ & 10 & 10 & 21 & $\bullet$ & $\bullet$ & $\bullet$ & $\bullet$ & $\bullet$ & $\bullet$ & $\bullet$ & $\bullet$ & 32 & 35 & $7^{10}$ & 50 & 200 & 2 &  &  &  &  &  &  &  &  \\
\bottomrule
\end{tabular}

\caption{The $32$ classes and the eight patterns.}\label{tab:patterns}
\end{table}

\end{document}
